\subsection{合取}

设 A、B 为组装。组装

\[
\text { 非 } ((\text { 非 } A) \text { 或 } (\text { 非 } B))
\]

将记为“A 且 B”。

CS6. 设 A、B、T 为组装，x 为一个字母。则组装

\[
(T | \boldsymbol {x}) (\boldsymbol {A} \text { 且 } \boldsymbol {B})
\]

与“(T\textbar x)A 和 (T\textbar x)B''。

这是CS5（§1，第2号）的直接推论。

CF9. 若A、B是C中的关系，则“A且B”是C中的关系（称为A与B的合取）。

这直接从CF1和CF2（§1，第4号）得出。

C20. 如果A、B是T中的定理，那么“A且B”是T中的定理。

假设“A且B”为假，也就是说，

\[
\text { 非   非   } ((\text { 非   } A) \text {   或   } (\text { 非   } B))
\]

是真的。根据C16，“（非A）或（非B）”，也就是说， \(A \Longrightarrow (\text{非 } B)\) 为真，因此“非B”为真；但这是荒谬的。因此“A且B”为真。

C21. 若A、B是C中的关系，则

\[
(A \text {   且   } B) \Longrightarrow A, \quad (A \text {   且   } B) \Longrightarrow B
\]

是 \(\mathcal{C}\) 中的定理 .

关系（非A） \(\Longrightarrow\) ((非A) 或 (非B)), (非B) \(\Longrightarrow\) ((非A)或(非B))在C中是定理，由S2（第1条）和C7（第2条）得出。现在((非A)或(非B)) \(\Longrightarrow\) (非(A且B)) 是 C 中由 C11 得到的一个定理。因此 (非A) \(\Longrightarrow\) (非(A且B)), (非B) \(\Longrightarrow\) (非(A且B)) 是 C 中的定理。结果通过应用 C17 得出。

我们用“A且B且C”（相应地，“A或B或C”）来表示关系“A且（B且C）”（相应地，“A或（B或C）”）。更一般地，如果

\[
A _ {1}, A _ {2}, \dots , A _ {n}
\]

是关系，我们用`` \(A_{1}\) 且 \(A_{2}\) ，以及 \ldots{} 且 \(A_{p}\) ''表示一种按如下约定逐步构造的关系，即`` \(A_{1}\) 且 \(A_{2}\) 且 \ldots{} 且 \(A_{h}\) ''表示与`` \(A_{1}\) 且 ( \(A_{2}\) 且 \ldots{} 且 \(A_{h}\) )''相同的关系。`` \(A_{1}\) 或 \(A_{2}\) 或 \ldots{} 或 \(A_{h}\) ''的关系类似地定义。`` \(A_{1}\) 且 \(A_{2}\) 且 \ldots{} 且 \(A_{h}\) ''是 \(\mathfrak{C}\) 中的定理，当且仅当每个关系 \(A_{1}, A_{2}, \ldots, A_{h}\) 都是 \(\mathfrak{C}\) 中的定理。

由此可见，每一个逻辑理论 \(\mathcal{C}\) 等价于一个至多有一个显式公理的逻辑理论 \(\mathcal{C}'\) 。这一点是显然的，如果 \(\mathcal{C}\) 没有显式公理。如果 \(\mathcal{C}\) 有显式公理 \(A_{1}, A_{2}, \ldots, A_{h}\) ，设 \(\mathcal{C}'\) 是这样一个理论，它具有与 \(\mathcal{C}\) 相同的符号和方案，以及显式公理`` \(A_{1}\) 且 \(A_{2}\) 且 \ldots{} 且 \(A_{h}\) ''。立即可见， \(\mathcal{C}\) （相应地， \(\mathcal{C}'\) ）的每一条公理都是 \(\mathcal{C}'\) （相应地， \(\mathcal{C}\) ）中的定理。

设 \(\mathfrak{C}_0\) 是这样一个理论：它没有显式公理，具有与 \(\mathfrak{C}\) 相同的符号，并将S1、S2、S3、S4作为其仅有的方案。那么对 \(\mathfrak{C}\) 的研究，原则上归结为对 \(\mathfrak{C}_0\) 的研究：对于关系 \(A\) 成为 \(\mathfrak{C}\) 中的定理，必要且充分的条件是存在公理 \(A_1, A_2, \ldots, A_h\)（\(\mathfrak{C}\) 的公理） ，使得 \((A_1\) 并且 \(A_2\) 并且 \(\ldots\) 并且 \(A_h) \Longrightarrow A\) 是 \(\mathfrak{C}_0\) 中的定理 。这个条件显然是充分的。反过来，假设 \(A\) 是 \(\mathfrak{C}\) 中的定理 ，并且设 \(A_1, A_2, \ldots, A_h\) 是 \(\mathfrak{C}\) 的那些公理，它们出现在 \(\mathfrak{C}\) 中一个包含 \(A\) 的证明中。设 \(\mathfrak{C}'\) （相应地， \(\mathfrak{C}''\) ）是由 \(\mathfrak{C}_0\) 通过附加公理 \(A_1, A_2, \ldots, A_h\) （相应地，公理`` \(A_1\) 并且 \(A_2\) 并且 \(\ldots\) 并且 \(A_h\) ''）而构造的理论。 \(A\) 在 \(\mathfrak{C}\) 中的证明也是 \(A\) 在 \(\mathfrak{C}'\) 中的一个证明，因此 \(A\) 是 \(\mathfrak{C}'\) 中的定理 ，从而也是 \(\mathfrak{C}''\) 中的定理，因为（如上所述） \(\mathfrak{C}'\) 和 \(\mathfrak{C}''\) 是等价的。根据演绎准则， \((A_1\) 并且 \(A_2\) 并且 \(\ldots\) 并且 \(A_h) \Longrightarrow A\) 是 \(\mathfrak{C}_0\) 中的定理 。如果 \(\mathfrak{C}\) 是矛盾的，那么由前述内容可知，存在一个合取式 \(A\)（由 \(\mathfrak{C}\) 的公理组成）以及一个关系 \(R\)（在 \(\mathfrak{C}\) 中），使得 \(A \Longrightarrow (R \text{ 且 } (\text{非 } R))\) 是 \(\mathfrak{C}_0\) 中的定理 。因此

\[
((\text { 非 } R) \text { 或 } (\text { 非   非 } R)) \Longrightarrow (\text { 非 } A)
\]

是 \(\mathfrak{G}_0\) 中的定理 ，并且由于``（非 \(R\) ）或（非非 \(R\) ）''是 \(\mathfrak{G}_0\) 中的一个定理，``非 \(A\) ''是 \(\mathfrak{G}_0\) 中的一个定理。反之，如果存在一个合取式 \(A\)（由 \(\mathfrak{G}\) 的公理组成），使得``非 \(A\) ''是 \(\mathfrak{G}_0\) 中的一个定理，那么 \(A\) 以及``非 \(A\) ''都是 \(\mathfrak{G}\) 中的定理，因此 \(\mathfrak{G}\) 是矛盾的。

